\documentclass[10pt]{article}
\usepackage[utf8]{inputenc}
\usepackage[vietnamese]{babel}
\usepackage{amsmath,amssymb}
\usepackage{geometry}
\geometry{
  paperwidth=8.5in,
  paperheight=11in,
  left=0.75in,
  right=0.75in,
  top=0.75in,
  bottom=0.75in,
  headheight=14pt,
  footskip=0.4in
}
\usepackage{multicol}
\usepackage{enumitem}
\usepackage{xcolor}
\usepackage{colortbl}
\usepackage{tikz}
\usepackage{fancyhdr}

\pagestyle{fancy}
\fancyhf{}
\renewcommand{\headrulewidth}{0pt}
\fancyhead[L]{\textsf{\textbf{SECTION 3.7}}\quad \textsf{Rates of Change in the Natural and Social Sciences}}
\fancyhead[R]{\textsf{\textbf{237}}}

\definecolor{tableblue}{RGB}{218, 234, 246}
\definecolor{techblue}{RGB}{0, 153, 204}

\newcommand{\techicon}{%
  \begin{tikzpicture}[baseline=(char.base)]
    \node[fill=techblue, text=white, inner sep=1.2pt, rounded corners=0.5pt, font=\sffamily\bfseries\tiny] (char) {T};
  \end{tikzpicture}%
}

\setlength{\columnsep}{22pt}
\setlength{\parindent}{0pt}
\setlength{\parskip}{4pt}

\begin{document}

\begin{multicols}{2}

\begin{enumerate}[label=\textbf{\arabic*.}, leftmargin=*, itemsep=10pt, topsep=0pt]
  \setcounter{enumi}{25}

  % Bài 26
  \item Nếu trong Ví dụ 4, một phân tử của sản phẩm C được tạo thành từ một phân tử chất phản ứng A và một phân tử chất phản ứng B, đồng thời nồng độ ban đầu của A và B có cùng giá trị $[\text{A}] = [\text{B}] = a\text{ mol/L}$, thì
  \[
  [\text{C}] = \frac{a^2kt}{akt + 1}
  \]
  trong đó $k$ là một hằng số.
  \begin{enumerate}[label=(\alph*), itemsep=2pt, topsep=2pt, leftmargin=*]
    \item Tìm tốc độ phản ứng tại thời điểm $t$.
    \item Chứng minh rằng nếu $x = [\text{C}]$, thì
    \[
    \frac{dx}{dt} = k(a - x)^2
    \]
    \item Điều gì xảy ra với nồng độ khi $t \to \infty$?
    \item Điều gì xảy ra với tốc độ phản ứng khi $t \to \infty$?
    \item Các kết quả ở phần (c) và (d) có ý nghĩa thực tế như thế nào?
  \end{enumerate}

  % Bài 27
  \item Trong Ví dụ 6 chúng ta đã xét một quần thể vi khuẩn nhân đôi sau mỗi giờ. Giả sử một quần thể vi khuẩn khác tăng gấp ba sau mỗi giờ và bắt đầu với 400 cá thể vi khuẩn. Hãy tìm biểu thức xác định số lượng $n$ vi khuẩn sau $t$ giờ và sử dụng nó để ước tính tốc độ tăng trưởng của quần thể vi khuẩn sau 2,5 giờ.

  % Bài 28
  \item Số lượng tế bào nấm men trong một môi trường nuôi cấy phòng thí nghiệm ban đầu tăng nhanh nhưng cuối cùng chững lại. Quần thể được mô hình hóa bởi hàm số
  \[
  n = f(t) = \frac{a}{1 + be^{-0{,}7t}}
  \]
  trong đó $t$ được đo bằng giờ. Tại thời điểm $t = 0$, quần thể có 20 tế bào và đang tăng với tốc độ 12 tế bào/giờ. Tìm các giá trị của $a$ và $b$. Theo mô hình này, điều gì sẽ xảy ra với quần thể nấm men về lâu dài?

  % Bài 29
  \item[\techicon~\textbf{29.}] Bảng số liệu cho biết dân số thế giới $P(t)$, tính theo đơn vị triệu người, trong đó $t$ được đo bằng năm và $t = 0$ tương ứng với năm 1900.

  \begin{center}
  \small
  \renewcommand{\arraystretch}{1.15}
  \begin{tabular}{|c|c|c|c|}
  \hline
  \rowcolor{tableblue}
  $t$ & \begin{tabular}{@{}c@{}} Dân số \\ (triệu) \end{tabular} & $t$ & \begin{tabular}{@{}c@{}} Dân số \\ (triệu) \end{tabular} \\
  \hline
  \phantom{0}0  & 1650 & \phantom{0}60  & 3040 \\
  10 & 1750 & \phantom{0}70  & 3710 \\
  20 & 1860 & \phantom{0}80  & 4450 \\
  30 & 2070 & \phantom{0}90  & 5280 \\
  40 & 2300 & 100 & 6080 \\
  50 & 2560 & 110 & 6870 \\
  \hline
  \end{tabular}
  \end{center}

  \begin{enumerate}[label=(\alph*), itemsep=3pt, topsep=2pt, leftmargin=*]
    \item Ước tính tốc độ tăng trưởng dân số vào năm 1920 và năm 1980 bằng cách lấy trung bình hệ số góc của hai cát tuyến.
    \item Sử dụng máy tính đồ họa hoặc máy vi tính để tìm một hàm bậc ba (đa thức bậc ba) mô hình hóa dữ liệu.
    \item Sử dụng mô hình ở phần (b) để tìm một mô hình cho tốc độ tăng trưởng dân số.
    \item Sử dụng phần (c) để ước tính tốc độ tăng trưởng vào năm 1920 và năm 1980. So sánh với các ước tính của bạn ở phần (a).
    \columnbreak
    \item Trong Mục 1.1, chúng ta đã mô hình hóa $P(t)$ bằng hàm mũ
    \[
    f(t) = (1{,}43653 \times 10^9) \cdot (1{,}01395)^t
    \]
    Sử dụng mô hình này để tìm một mô hình cho tốc độ tăng trưởng dân số.
    \item Sử dụng mô hình ở phần (e) để ước tính tốc độ tăng trưởng vào năm 1920 và năm 1980. So sánh với các ước tính của bạn ở phần (a) và (d).
    \item Ước tính tốc độ tăng trưởng vào năm 1985.
  \end{enumerate}

  % Bài 30
  \item[\techicon~\textbf{30.}] Bảng số liệu cho thấy độ tuổi kết hôn trung bình lần đầu của phụ nữ Nhật Bản đã biến thiên như thế nào kể từ năm 1950.

  \begin{center}
  \small
  \renewcommand{\arraystretch}{1.15}
  \begin{tabular}{|c|c|c|c|}
  \hline
  \rowcolor{tableblue}
  $t$ & $A(t)$ & $t$ & $A(t)$ \\
  \hline
  1950 & 23{,}0 & 1985 & 25{,}5 \\
  1955 & 23{,}8 & 1990 & 25{,}9 \\
  1960 & 24{,}4 & 1995 & 26{,}3 \\
  1965 & 24{,}5 & 2000 & 27{,}0 \\
  1970 & 24{,}2 & 2005 & 28{,}0 \\
  1975 & 24{,}7 & 2010 & 28{,}8 \\
  1980 & 25{,}2 & 2015 & 29{,}4 \\
  \hline
  \end{tabular}
  \end{center}

  \begin{enumerate}[label=(\alph*), itemsep=3pt, topsep=2pt, leftmargin=*]
    \item Sử dụng máy tính đồ họa hoặc máy vi tính để mô hình hóa các dữ liệu này bằng một đa thức bậc bốn.
    \item Sử dụng phần (a) để tìm một mô hình cho $A'(t)$.
    \item Ước tính tốc độ thay đổi độ tuổi kết hôn của phụ nữ vào năm 1990.
    \item Vẽ đồ thị các điểm dữ liệu và các mô hình cho $A$ và $A'$.
  \end{enumerate}

  % Bài 31
  \item Tham khảo định luật dòng chảy tầng đã cho trong Ví dụ 7. Xét một mạch máu có bán kính $0{,}01\text{ cm}$, chiều dài $3\text{ cm}$, độ chênh áp $3000\text{ dyn/cm}^2$, và độ nhớt $\eta = 0{,}027$.
  \begin{enumerate}[label=(\alph*), itemsep=2pt, topsep=2pt, leftmargin=*]
    \item Tìm vận tốc của dòng máu dọc theo đường tâm $r = 0$, tại bán kính $r = 0{,}005\text{ cm}$, và tại thành mạch $r = R = 0{,}01\text{ cm}$.
    \item Tìm gradient vận tốc tại $r = 0$, $r = 0{,}005$, và $r = 0{,}01$.
    \item Vận tốc đạt giá trị lớn nhất ở đâu? Vận tốc biến thiên nhanh nhất ở đâu?
  \end{enumerate}

  % Bài 32
  \item Tần số dao động của một dây vĩ cầm (violin) đang rung được cho bởi công thức
  \[
  f = \frac{1}{2L} \sqrt{\frac{T}{\rho}}
  \]
  trong đó $L$ là chiều dài của dây, $T$ là lực căng, và $\rho$ là mật độ dài của nó. [Xem Chương 11 trong cuốn \textit{Musical Acoustics} của D. E. Hall, tái bản lần 3 (Pacific Grove, CA, 2002).]
  \begin{enumerate}[label=(\alph*), itemsep=2pt, topsep=2pt, leftmargin=*]
    \item Tìm tốc độ thay đổi của tần số theo
    \begin{enumerate}[label=(\roman*), itemsep=1pt, topsep=1pt, leftmargin=*]
      \item chiều dài (khi $T$ và $\rho$ không đổi),
      \item lực căng (khi $L$ và $\rho$ không đổi), và
      \item mật độ dài (khi $L$ và $T$ không đổi).
    \end{enumerate}
    \item Cao độ của một nốt nhạc (độ bổng hay trầm của âm thanh phát ra) được quyết định bởi tần số $f$. (Tần số càng cao, cao độ càng bổng.) Hãy sử dụng dấu của các
  \end{enumerate}

\end{enumerate}

\end{multicols}

\vfill
\begin{center}
\scriptsize\color{black!70}
Copyright 2021 Cengage Learning. All Rights Reserved. May not be copied, scanned, or duplicated, in whole or in part. WCN 02-200-203
\end{center}

\end{document}